\uuid{Hwoa}
\titre{ MSE, MAE et dépendance à la taille du jeu de données }
\niveau{L3}
\module{Analyse de données}
\chapitre{Réseaux de neurones}
\sousChapitre{Fonctions de coût, régression}
\theme{MSE, MAE, comparaison de modèles, taille du dataset}
\auteur{Maxime Nguyen}
\datecreate{2026-03-19}
\organisation{AMSCC}
\difficulte{2}

\contenu{
\texte{
  On compare trois implémentations d'une fonction de coût :
  \begin{itemize}
    \item[A.] \texttt{J = np.mean((y\_hat - y) ** 2)}
    \item[B.] \texttt{J = np.sum((y\_hat - y) ** 2)}
    \item[C.] \texttt{J = np.mean(np.abs(y\_hat - y))}
  \end{itemize}
}
\begin{enumerate}

\item
  \question{Laquelle est la MSE standard ? Laquelle est la MAE ?}
  \reponse{
    \textbf{A} est la MSE standard (moyenne des carrés des erreurs).
    \textbf{C} est la MAE (moyenne des valeurs absolues des erreurs).
  }

\item
  \question{Pourquoi B est-il problématique si on compare deux modèles entraînés sur des jeux de données de tailles différentes ?}
  \reponse{
    B calcule la \emph{somme} (et non la moyenne) des erreurs au carré.
    Sa valeur dépend directement du nombre d'exemples $m$ : un modèle entraîné sur $m = 10000$ exemples
    aura mécaniquement un coût beaucoup plus élevé qu'un modèle entraîné sur $m = 100$,
    même si les erreurs individuelles sont identiques. La comparaison est donc impossible.
  }

\end{enumerate}
}
